\section{Boundaries}

Forensic analytics can prioritize investigation. They cannot prove fraud,
certify correctness, or replace risk-limiting audits. Residual methods require
that sentence because they produce familiar-looking statistics whose limits are
easy to forget.

The main technical boundary is baseline quality. A sparse, biased, stale, or
misaligned baseline can manufacture false outliers. If a precinct changed shape,
if comparable units are chosen poorly, if demographic covariates omit a new
population center, or if historical turnout came from a different contest type,
then the residual is partly measuring the baseline defect. The report must name
the feature set and baseline population so that this weakness is visible.

The second boundary is causal. A residual is descriptive: it says that an
observed turnout rate or vote share was far from what the model predicted. It
does not say why. The explanation may be administrative, demographic,
political, geographic, or data-quality related. Human investigation and, where
appropriate, audit procedures remain outside the residual statistic.
